\section{Free compression and the Hankel test}
\label{sec:hankel}

Retain $e,p,d,w,u$, the isometries $U_k$, and the subspace $\mathcal K$
from the preceding section. We construct $h\in e\cL e$ whose norm is
of order $u$ and whose compression to $\mathcal K$ has a Hankel component
of the same order, with a uniform lower bound after division by $u$.

\subsection{The dihedral corner}

Let $\mathcal C=W^*(d,s)$. Freeness of the centered symmetries $d,s$
shows that their alternating reduced words, together with $1$, form an
orthonormal word basis of $L^2(\mathcal C)$. These are precisely the
group words of the infinite dihedral group. In particular,
\[
  w=ds\text{ is a Haar unitary},\qquad dwd=w^{-1},\qquad p=(1+d)/2.
\]
Write $m$ for Haar probability measure on $\T$.

\begin{lemma}[The abelian corner and corner freeness]
\label{lem:dihedral-corner}
The algebra $p\mathcal C p$ is identified with the even functions on the
circle by
\begin{equation}
  \begin{gathered}
    p\mathcal C p
      =\{f(w)p:f\in L^\infty(\T,m),\ f(z)=f(z^{-1})\},\\
    \tau_p(f(w)p)=\int_{\T}f\,dm.
  \end{gathered}
  \label{eq:dihedral-corner}
\end{equation}
The map $f\mapsto f(w)p$ in this formula is a faithful normal unital
$*$-isomorphism onto the corner. Moreover, $pDp$ and $p\mathcal C p$
are free with respect to the normalized trace $\tau_p$.
\end{lemma}

\begin{proof}
Every group word has the form $w^l$ or $w^ld$, $l\in\mathbb Z$.
Since $dp=p$ and $dw^ld=w^{-l}$,
\[
  pw^ldp=pw^lp=\frac{w^l+w^{-l}}2p.
\]
An even function of $w$ commutes with $d$ and hence with $p$. Thus
$f\mapsto f(w)p$ is a normal $*$-homomorphism on the even-function
algebra, with unit $p$. The group trace gives
\[
  \tau_p(f(w)p)
      =2\tau(f(w)p)=\tau(f(w))+\tau(f(w)d)=\int_{\T}f\,dm,
\]
initially for trigonometric polynomials and then for bounded functions
by $L^2$ approximation. In particular the map is faithful, and its
range is a von Neumann algebra. The compressed group-word span is
strongly dense in $p\mathcal C p$ and consists of even Laurent
polynomials times $p$. This proves \eqref{eq:dihedral-corner} and the
isomorphism assertion. It also gives
\[
  \|f(w)p\|=\|f\|_\infty,\qquad
  \|f(w)p\|_{2,\tau_p}=\|f\|_{L^2(m)}.
\]

For freeness, let $a_1,\ldots,a_j\in pDp$ be centered for $\tau_p$.
Then $\tau(a_i)=0$ and $a_id=a_i$. The centered trigonometric span in
$p\mathcal C p$ is generated by the elements $pw^lp$, $l\ge1$.
For $l_1,\ldots,l_j\ge1$, absorbing the corner projections and using
traciality gives
\begin{align*}
 &\tau_p(a_1pw^{l_1}p\cdots a_jpw^{l_j}p)\\
 &\qquad=2\tau(a_1w^{l_1}\cdots a_jw^{l_j})\\
 &\qquad=2\tau\big(a_1s(ds)^{l_1-1}\cdots
                              a_js(ds)^{l_j-1}\big)=0.
\end{align*}
All the displayed $D$ letters are centered, and $s$ is centered and
free from the whole algebra $D$, by $\cL=D*A_0$. Thus the last word is
a reduced alternating word of trace zero. Fej\'er means of an even
bounded zero-mean function are even zero-mean trigonometric polynomials,
are uniformly bounded, and converge in $L^2(m)$. Approximating each
corner letter in this manner extends the vanishing to all cyclic
alternating products of even length. Products of odd alternating
length have their endpoints in the same algebra. Move one endpoint
cyclically next to the other, split their product into its centered and
scalar parts, and apply the even-length case and induction on the
length. The length-one case is centeredness. This proves freeness.
\end{proof}

\subsection{The expectation and the norm of a free compression}

\begin{lemma}[Free compression]
\label{lem:free-compression}
For $b=b^*\in p\mathcal C p$ with $\tau_p(b)=0$,
\begin{equation}
  E_{\mathcal C}(ebe)=u^2b,\qquad
  \|ebe\|\le2\sqrt u\,\|b\|_{2,\tau_p}+u\|b\|,
  \qquad u=\tau_p(e).
  \label{eq:free-compression}
\end{equation}
Here $E_{\mathcal C}$ is the trace-preserving conditional expectation
from $\cL$ onto $\mathcal C$.
\end{lemma}

\begin{proof}
Put $\mathcal A=pDp$, $\mathcal Z=p\mathcal C p$, and
$a=e-up\in\mathcal A$, so $\tau_p(a)=0$. Since $p\in\mathcal C$,
bimodularity shows that $E_{\mathcal C}(ebe)$ belongs to $\mathcal Z$.
On $p\cL p$, its restriction is the trace-preserving expectation onto
$\mathcal Z$ for $\tau_p$. Expand
\[
  ebe=u^2b+u(ab+ba)+aba.
\]
Freeness gives zero expectation for $ab$ and $ba$. For the last term,
test against any $c\in\mathcal Z$ and write
$c=c^\circ+\tau_p(c)p$. Freeness and traciality give
\[
  \tau_p(c^\circ aba)=0,
  \qquad \tau_p(aba)=\tau_p(ba^2)
                  =\tau_p(b)\tau_p(a^2)=0.
\]
Thus $E_{\mathcal Z}(aba)=0$ also. This proves the expectation formula.

For the norm estimate work in
\[
  H_c=L^2(W^*(\mathcal A,\mathcal Z),\tau_p).
\]
This representation is faithful, so its operator norm is the algebra
norm. The free-word decomposition, grouping a possible initial
$\mathcal A$ letter, takes the form
\[
  H_c=L^2(\mathcal A,\tau_p)\otimes F_0,
  \qquad F_0=\C\Omega\oplus F_0^\circ,
\]
where $F_0^\circ$ is the closed span of reduced words starting in the
centered $L^2$ subspace $\mathcal Z^\circ$. Let $Q_0$
be the projection of $F_0$ onto the complement of the vacuum, and let
$\Pi$ project in $H_c$ onto words starting in $\mathcal Z^\circ$.
Then
\[
  \Pi=|1_{\mathcal A}\rangle\langle1_{\mathcal A}|\otimes Q_0,
  \qquad L_e=\lambda_{\mathcal A}(e)\otimes I_{F_0},
\]
where $1_{\mathcal A}=p$ has $L^2(\tau_p)$ norm one. Consequently
\begin{equation}
  \|L_e\Pi\|=\|\Pi L_e\|\le\|e\|_{2,\tau_p}=\sqrt u.
  \label{eq:corner-overlap}
\end{equation}

The following three blocks are the standard creation, annihilation,
and diagonal decomposition of a centered free-product letter; see
\citet[Section~1, Lemmas~2.1--2.3 and Corollary~2.4]{RicardXu2006}.
We give the bounds needed here directly.
Let $T=L_b$ on $H_c$. Because $b$ is centered, on the vacuum and on
words starting in $\mathcal A^\circ$ its left action simply prepends
the letter $b$. Therefore
\[
  (1-\Pi)T(1-\Pi)=0,\qquad
  \|\Pi T(1-\Pi)\|\le\|b\|_{2,\tau_p}.
\]
Self-adjointness gives the same bound for $(1-\Pi)T\Pi$, and
$\|\Pi T\Pi\|\le\|b\|$. Decomposing $T$ into these three blocks,
using $\|L_e\|\le1$, and applying \eqref{eq:corner-overlap}, we obtain
\begin{align*}
  \|ebe\|=\|L_eTL_e\|
   &\le\|L_e\Pi T(1-\Pi)L_e\|
      +\|L_e(1-\Pi)T\Pi L_e\|
      +\|L_e\Pi T\Pi L_e\|\\
   &\le\sqrt u\,\|b\|_{2,\tau_p}
      +\sqrt u\,\|b\|_{2,\tau_p}+u\|b\|.
\end{align*}
This proves \eqref{eq:free-compression}.
\end{proof}

\subsection{Fourier coefficients of the compressed operator}

The preceding lemma gives both the expectation and the upper norm
bound for a corner compression. We now compute its matrix entries on
$\mathcal K$; their dependence on the sum of the indices will isolate
the commutator with $L_s$.

Let $f\in L^\infty(\T,m)$ be real, even, and of mean zero, and set
\[
  b_f=f(w)p,\qquad h=eb_fe,\qquad
  \gamma_l=\int_{\T}z^lf(z)\,dm(z)\quad(l\in\mathbb Z).
\]
Then $b_f$ meets the hypotheses of Lemma~\ref{lem:free-compression}.
We claim that, for $j,k\ge0$,
\begin{equation}
  \begin{aligned}
    U_j^*L_hU_k
      &=L_{t^{-1}E_{A_0}(w^jhw^{-k})}\big|_B\\
      &=u\big(\gamma_{j-k}I_B+\gamma_{j+k+1}L_s|_B\big).
  \end{aligned}
  \label{eq:hankel-compression}
\end{equation}
For the first equality, test on $v,z\in A_0$. Since $eh=he=h$,
\begin{align*}
  \langle L_hU_kv,U_jz\rangle
    &=t^{-1}\tau(z^*w^jhw^{-k}v)\\
    &=\left\langle
        t^{-1}E_{A_0}(w^jhw^{-k})v,z\right\rangle.
\end{align*}
Density then gives the operator identity on $B$.

The reduced-word decomposition shows that on $W^*(D,s)$ the expectation
$E_{A_0}$ equals the expectation onto $W^*(s)=\C1+\C s$: every reduced
word containing a centered $D$ letter is orthogonal to $L^2(A_0)$.
Because $W^*(s)\subset\mathcal C$, this latter expectation factors
through $E_{\mathcal C}$. By \eqref{eq:free-compression} and
bimodularity,
\[
  E_{\mathcal C}(w^jhw^{-k})=u^2w^jb_fw^{-k}.
\]
Using $dw^{-k}=w^kd$ and $d=ws$, compute
\[
  w^jpw^{-k}
     =\frac{w^{j-k}+w^{j+k}d}{2}
     =\frac{w^{j-k}+w^{j+k+1}s}{2}.
\]
The group trace implies
\[
  E_{W^*(s)}(f(w)w^l)=\gamma_l1,\qquad
  E_{W^*(s)}(f(w)w^ls)=\gamma_ls.
\]
For instance this follows on Laurent polynomials because only the
identity power contributes, and then by $L^2$ approximation. Hence
\[
  t^{-1}E_{A_0}(w^jhw^{-k})
     =\frac{u^2}{2t}
          (\gamma_{j-k}1+\gamma_{j+k+1}s)
     =u(\gamma_{j-k}1+\gamma_{j+k+1}s),
\]
which proves \eqref{eq:hankel-compression}, including its normalization
and its Hankel index.

The matrices
\[
  T_f=(\gamma_{j-k})_{j,k\ge0},\qquad
  H_f=(\gamma_{j+k+1})_{j,k\ge0}
\]
define bounded operators on $\ell^2(\N_0)$ of norm at most
$\|f\|_\infty$. To see this directly, decompose
\[
  L^2(\T,m)=H_-\oplus H_+,\qquad
  H_-=\overline{\operatorname{span}}\{z^{-j}:j\ge0\},\quad
  H_+=\overline{\operatorname{span}}\{z^{k+1}:k\ge0\}.
\]
In the indicated orthonormal bases, $T_f$ and $H_f$ are respectively
the $H_-,H_-$ and $H_-,H_+$ blocks of multiplication by $f$.
Equation~\eqref{eq:hankel-compression}, first on finite coordinate
vectors and then by continuity, therefore gives
\begin{equation}
  \iota^*L_h\iota
      =u(T_f\otimes I_B+H_f\otimes L_s|_B).
  \label{eq:toeplitz-hankel-form}
\end{equation}
As $J=I\otimes B_0$, it follows that
\begin{equation}
  \big\|[J,\iota^*L_h\iota]\big\|
       =u\|H_f\|\,\|[B_0,L_s|_B]\|.
  \label{eq:hankel-commutator}
\end{equation}
Here we used the operator norm identity
$\|A\otimes B\|=\|A\|\|B\|$ for Hilbert-space tensor products.

\begin{proof}[Completion of the proof of Theorem~\ref{thm:free-corner}]
Put $\theta=u/4=1/(2N)$, so $0<\theta\le1/4$, and use angular coordinates
$z=e^{i\alpha}$, $-\pi<\alpha\le\pi$, to define
\[
  f(e^{i\alpha})=\mathbf1_{[-\theta,\theta]}(\alpha)-\frac\theta\pi.
\]
This function is real, even, and of mean zero. The trace formula in
\eqref{eq:dihedral-corner} gives
\[
  \|b_f\|\le1,\qquad
  \|b_f\|_{2,\tau_p}^2
       =\frac\theta\pi\left(1-\frac\theta\pi\right)\le u.
\]
Consequently \eqref{eq:free-compression} gives
\begin{equation}
  \|h\|=\|eb_fe\|\le3u.
  \label{eq:narrow-arc-norm}
\end{equation}
For every $l\ge1$ the Fourier coefficient is
\[
  \gamma_l=\frac{\sin(l\theta)}{\pi l}.
\]
If $0\le j,k<N$, then $1\le j+k+1\le2N-1$ and
$(j+k+1)\theta<1<\pi/2$. The inequality
$\sin x\ge2x/\pi$ on $[0,\pi/2]$ therefore gives
\[
  (H_f)_{jk}\ge\frac{2\theta}{\pi^2}.
\]
Testing the leading $N$ by $N$ block on its constant unit vector,
and using $\theta N=1/2$, proves the uniform bound
\begin{equation}
  \|H_f\|\ge\frac{2\theta N}{\pi^2}\ge\frac1{2\pi^2}.
  \label{eq:narrow-arc-hankel}
\end{equation}
Combine \eqref{eq:hankel-commutator}, \eqref{eq:compression-bound},
\eqref{eq:narrow-arc-norm}, and \eqref{eq:narrow-arc-hankel} to obtain
\[
  \frac{u}{2\pi^2}\|[B_0,L_s|_B]\|
       \le3u(1+4c_0)\epsilon.
\]
Cancelling $u>0$ yields
\[
  \|[B_0,L_s|_B]\|
      \le6\pi^2(1+4c_0)\epsilon=c_1\epsilon.
\]
The level-zero case of \eqref{eq:repeated-letter} identifies $G|_B$
with $B_0$, so this is exactly the required compression estimate for
the symmetrized operator. Lemma~\ref{lem:free-symmetrization} preserves
the compression of $[G,L_s]$ between vectors of $B$. The same estimate
therefore holds for the original operator in
Theorem~\ref{thm:free-corner}, completing its proof.
\end{proof}
